% !TeX spellcheck = fr-FR, en-US
% !TeX encoding = UTF-8
% !TEX program = pdflatex

% VSCODE word wrap: ALT + Z
% COMPILE WITH:
% `latexmk`
% latexmk -pdf main.tex
% You need pdflatex and biber (in all TeXLive distributions)

\documentclass[11pt]{article} % text width
\usepackage[utf8]{inputenc} % encode text to utf8

% Set the label for itemize globally, if desired
\usepackage{enumitem}
\setlist[itemize]{label=\textbullet}

% paragraph formatting: https://www.overleaf.com/learn/latex/Paragraph_formatting
\setlength{\parindent}{1em}
\setlength{\parskip}{1em}

% better language support
\usepackage[english, french]{babel}

% use pdflatex
\usepackage[T1]{fontenc} % font encoding
\usepackage[a4paper, margin=2cm, head=18.0pt]{geometry} % set margins to 1.5 cm
\usepackage{graphicx}% for graphics
\usepackage[onehalfspacing]{setspace}
\usepackage{tocbasic}
\usepackage{booktabs}
\usepackage{multicol}
\usepackage{multirow}
\usepackage[]{scrlayer-scrpage}
\usepackage[titletoc]{appendix}
\usepackage{comment}

\usepackage{dirtree} % directory tree
\usepackage{float} % float images
\usepackage{microtype} % better text formatting, avoid one word on a new line for example
\usepackage{mathtools} % math tools
\usepackage{amsmath, amssymb} % math symbols
\usepackage{algorithm} % algorithm
% \usepackage{algorithmic}  % old, use algorithmicx from algpseudocode
\usepackage{algpseudocode} % pseudocode

% fix Font shape `T1/fve/m/sc' undefined on \Procedure
\renewcommand{\algorithmicprocedure}{\textbf{Procedure}}

\usepackage{tikz} % for drawing graphs
\usetikzlibrary{arrows.meta} % for arrow styles
%\usetikzlibrary{positioning} % for relative positioning of nodes

% quotes and bibliography: https://www.overleaf.com/learn/latex/Typesetting_quotations
\usepackage{csquotes}
\usepackage[
    % Use the french style of quotes, which are more visibly distinct
    left = \flqq{},% 
    right = \frqq{},% 
    leftsub = \flq{},% 
    rightsub = \frq{} %
]{dirtytalk}
\DeclareQuoteStyle{english}{\glqq}{\grqq}{\glq}{\grq}

% \usepackage[
%     backend=biber,
%     style=numeric,
%     sorting=none
% ]{biblatex}
%\usepackage[backend=biber, style=numeric, defernumbers=true, language=american]{biblatex}
\usepackage[backend=biber, style=numeric, sorting=none, defernumbers=true, language=american]{biblatex}
% add commands for automatic cite/uncite distinction
\DeclareBibliographyCategory{cited}
\AtEveryCitekey{\addtocategory{cited}{\thefield{entrykey}}}
\addbibresource{biblio.bib} % bibliography
\nocite{*} % all references

\newcommand{\ts}{\textsuperscript} % superscript for 2nd or XIXème

\pagenumbering{roman} % set page numbering of front matter sections

% use acronyms and glossaries
% toc: add glossary to table of contents
\usepackage{hyperref}
\usepackage{xurl}
\usepackage[acronym, toc]{glossaries} 
\makeglossaries
\newglossaryentry{memgraph}
{
  name=memory graph,
  description={A memory graph, or \textit{memgraph} is a graph representation of a memory dump. This graph can be a graph of blocks, where each node in the graph corresponds to a block of 8 bytes in the heap dump and each edge corresponds to a pointer from one block to another, or describes which blocks are part of a chunk whose root note is a Chunk Header Node. It can also be a graph of chunks (only CHNs), where each node in the graph corresponds to a chunk in heap dump and each edge corresponds to a pointer from one object to another.}
}

\newglossaryentry{nodes}
{
  name=nodes,
  description={A node is an entity in a graph, it can be a person, a place, a thing, or any other entity.}
}
\newglossaryentry{chn}{
  name={CHN},
  description={Chunk Header Node. This is a node whose bytes have been identified as a data structure header. In the graph, this node is the root node of an malloc-allocated memory chunk.}
}

\newglossaryentry{pn}{
  name={PN},
  description={Pointer Node. This is a node whose bytes have been identified as a pointer.}
}

\newglossaryentry{kn}{
  name={KN},
  description={Key Node. This is a node whose bytes have been identified as a key. This identification relies both on the annotations and some verification checks.}
}

\newglossaryentry{vn}{
  name={VN},
  description={Value Node. These are all blocks that have not been identified. It is the default node type.}
}

\newacronym{kg}{KG}{Knowledge Graph}
\newacronym{foss}{FOSS}{Free and Open Source Software}
\newacronym{rdf}{RDF}{Resource Description Framework}
\newacronym{rdfs}{RDFS}{Resource Description Framework Schema}
\newacronym{owl}{OWL}{Web Ontology Language}
\newacronym{ml}{ML}{Machine Learning}
\newacronym{dl}{DL}{Deep Learning}
\newacronym{fe}{FE}{Feature Evaluation}
\newacronym{nlp}{NLP}{Natural Language Processing}
\newacronym{ke}{KE}{Knowledge Engineering}
\newacronym{del}{DEL}{Directed Edge-labelled Graphs}
\newacronym{er}{ER}{Entity Resolution}
\newacronym{qa}{QA}{Quality Assurance}
\newacronym{sparql}{SPARQL}{SPARQL Protocol and RDF Query Language}
\newacronym{ssh}{SSH}{Secure Shell Protocol}
\newacronym{os}{OS}{Operating System}
\newacronym{vm}{VM}{Virtual Machine}
\newacronym{ddos}{DDoS}{Distributed Denial of Service}
\newacronym{ess}{ESS}{Estimated Security Strength}
\newacronym{vmi}{VMI}{Virtual Machine Introspection}
\newacronym{smote}{SMOTE}{Synthetic Minority Over-sampling Technique}
\newacronym{svm}{SVM}{Support Vector Machine}
\newacronym{knn}{KNN}{K-Nearest Neighbors}
\newacronym{rf}{RF}{Random Forest}
\newacronym{mlp}{MLP}{Multi-Layer Perceptron}
\newacronym{relu}{ReLU}{Rectified Linear Unit}
\newacronym{sgd}{SGD}{Stochastic Gradient Descent}
\newacronym{ai}{AI}{Artificial Intelligence}
\newacronym{pca}{PCA}{Principal Component Analysis}
\newacronym{lda}{LDA}{Linear Discriminant Analysis}
\newacronym{tsne}{t-SNE}{t-distributed Stochastic Neighbor Embedding}
\newacronym{msb}{MSB}{Most Significant Bit}
\newacronym{lsb}{LSB}{Least Significant Bit}

%deep learning
\newacronym{lstm}{LSTM}{Long Short-Term Memory}
\newacronym{gru}{GRU}{Gated Recurrent Units}
\newacronym{rnn}{RNN}{Recurrent Neural Networks}
\newacronym{cnn}{CNN}{Convolutional Neural Networks}
\newacronym{rcnn}{RCNN}{Recurrent Convolutional Neural Network}
\newacronym{gnn}{GNN}{Graph Neural Network}
\newacronym{gcn}{GCN}{Graph Convolutional Networks}
\newacronym{llm}{LLM}{Large Language Model}

% INSA
\newacronym{insa}{INSA}{Institut National des Sciences Appliquées}
\newacronym{aaia}{AAIA}{Algorithmique Avancée pour l'Intelligence Artificielle et les graphes}

%\input{glossaries.tex} % acronyms definitions, failed to make in work on a separate file!!!

% custom commands
% escape char in latex: https://tex.stackexchange.com/questions/34580/escape-character-in-latex
% horizontal spacing: https://tex.stackexchange.com/questions/74353/what-commands-are-there-for-horizontal-spacing/74354
\newcommand{\p}{\texttt{+}} % small unary plus
\newcommand{\doublep}{\texttt{++}} % double small unary plus
\newcommand{\m}{\texttt{-} \space} % small unary minus
\newcommand{\doublem}{\texttt{-}\texttt{-} \space} % double small unary minus

% title and section formatting
\usepackage{titlesec}

\setcounter{tocdepth}{3} % set depth of table of contents
\setcounter{secnumdepth}{3}  % Numbering depth of sections

% ------------------------------ code ------------------------------
\usepackage{listings} % code listings

% code listing style
\usepackage{bera}% optional: just to have a nice mono-spaced font
\usepackage{listings}
\usepackage{xcolor}

\definecolor{eclipseStrings}{RGB}{42,0.0,255}
\definecolor{eclipseKeywords}{RGB}{127,0,85}
\definecolor{punctuationcolor}{rgb}{0.5,0,0}
\definecolor{delimcolor}{rgb}{0,0.5,0}
\definecolor{red}{rgb}{1,0,0}
\colorlet{numb}{magenta!60!black}

\lstset{
  language=bash,
  basicstyle=\ttfamily\small,
  breaklines=true,
  frame=single,
  numbers=left,
  numberstyle=\tiny,
  showstringspaces=false,
  tabsize=4,
  captionpos=b
}

\lstdefinestyle{json}{
  basicstyle=\ttfamily\small,
  breaklines=true,
  postbreak=\mbox{\space},
  columns=fullflexible,
  showstringspaces=false,
  commentstyle=\color{gray},
  keywordstyle=\color{black},
  numberstyle=\tiny\color{gray},
  numbers=left,
  frame=single,
  captionpos=b
}

\lstdefinestyle{text}{
  basicstyle=\ttfamily\small,
  breaklines=true,
  postbreak=\mbox{\space},
  columns=fullflexible,
  showstringspaces=false,
  commentstyle=\color{gray},
  keywordstyle=\color{black},
  numberstyle=\tiny\color{gray},
  numbers=left,
  frame=single,
  captionpos=b
}

\lstdefinestyle{hexdump}{
  basicstyle=\ttfamily\small,
  breaklines=true,
  postbreak=\mbox{\space},
  columns=fullflexible,
  showstringspaces=false,
  commentstyle=\color{gray},
  keywordstyle=\color{black},
  numberstyle=\tiny\color{gray},
  numbers=left,
  frame=single,
  captionpos=b
}

\lstdefinestyle{rust}{
  basicstyle=\ttfamily\small,
  breaklines=true,
  postbreak=\mbox{\space},
  columns=fullflexible,
  showstringspaces=false,
  commentstyle=\color{gray},
  keywordstyle=\color{black},
  numberstyle=\tiny\color{gray},
  numbers=left,
  frame=single,
  captionpos=b
}

\usepackage{hyphenat} % fix "overfull hbox" with slicing words using hyphenation
\hyphenation{hy-phen-a-tion} % indicate all 3 permissible hyphenation points

% where to put all images and figures
\graphicspath{{img/}}

% customize the header and footer of the document
\usepackage{scrlayer-scrpage}
\clearpairofpagestyles
\cfoot[\pagemark]{\pagemark}

% !TeX spellcheck = fr_FR, en_US
% !TeX encoding = UTF-8

% document info
\newcommand{\thetitle}{AAIA TP 3IF - Monde des Blocs}
\newcommand{\theauthor}{Florian Rascoussier \and Christine Solnon}

\title{\thetitle}
\author{\theauthor}
\date{\today} % title information

% document content
\begin{document}

% -- Title
\maketitle
% WARN: It's not really possible to add an image before the title, without a real title page
{ % WARN: The centering command needs its own scope
  \centering
  \includegraphics[width=7cm, height=1.5cm]{Logo_INSA.png}\\
}
%%%%%%%%%%%%%%%%%%%%%%%%%%%%%%%%%%%%%%%%%%%%%%%%%%%%%%%%%%%%%%%%%%%%%%%%%%%%%%%%%%%%%%%%%

% -- Abstract, Acknowledgements, Introduction...
% !TeX spellcheck = fr
% !TeX encoding = UTF-8

% -- Introduction
\section*{Introduction}
Le problème du monde des blocs est un cas classique en intelligence artificielle et en planification automatique. Il consiste à manipuler un ensemble de blocs disposés sur différentes piles pour atteindre une configuration spécifique.

Dans sa modélisation standard, le monde des blocs est représenté comme un problème de planification, où chaque configuration de blocs représente un état. Le but est de trouver une séquence d'actions (mouvements de blocs d'une pile à une autre) permettant de passer de l'état initial à l'état final désiré. Ceci est typiquement représenté sous la forme d'un graphe d'états, où chaque nœud correspond à une configuration de blocs et chaque arête à un mouvement possible entre deux configurations.

Les algorithmes A* et AWA* sont utilisés pour résoudre ce problème en cherchant le chemin le plus court dans le graphe d'états. A* est un algorithme de recherche informé qui utilise des heuristiques pour guider la recherche vers l'état final de manière efficace. AWA*, quant à lui, est une variante d'A* permettant une recherche anytime, adaptative en fonction des ressources disponibles.

Ce travail pratique (TP) se concentre sur l'application et la comparaison de ces algorithmes dans le contexte du monde des blocs. Nous discuterons en détail la mise en œuvre de ces algorithmes, l'importance de choisir des heuristiques adéquates, et comment elles peuvent drastiquement influencer les performances et l'efficacité de la recherche. L'objectif est de comprendre non seulement comment fonctionnent ces algorithmes de planification, mais aussi d'apprécier l'impact significatif des heuristiques pour guider la recherche vers des solutions optimales ou satisfaisantes dans des délais raisonnables.

En somme, ce TP offre une exploration approfondie du monde des blocs comme un microcosme pour comprendre les défis et les stratégies de la planification et de l'optimisation dans les systèmes intelligents.


%%%%%%%%%%%%%%%%%%%%%%%%%%%%%%%%%%%%%%%%%%%%%%%%%%%%%%%%%%%%%%%%%%%%%%%%%%%%%%%%%%%%%%%%%

\newpage

% francisation des noms de tableaux, figures, etc.
\renewcommand{\contentsname}{Table des Matières}
\def\lstlistingautorefname{Cd.}
\renewcommand{\lstlistingname}{Code}
\renewcommand{\lstlistlistingname}{Codes et programmes}

% -- Table of contents
\tableofcontents
\listoffigures
% \listoftables
\lstlistoflistings
% \newpage

%%%%%%%%%%%%%%%%%%%%%%%%%%%%%%%%%%%%%%%%%%%%%%%%%%%%%%%%%%%%%%%%%%%%%%%%%%%%%%%%%%%%%%%%%
\pagenumbering{arabic} % reset page numbering of main matter sections

\newpage

% -- Content
% !TeX spellcheck = fr
% !TeX encoding = UTF-8

% -- Exercice 1
\section{Partie 1 : Modélisation du problème}

Dans cette partie, on va s'intéresser à la modélisation du problème du \textit{Monde des Blocs} sous forme d'un problème de planification.

\subsection{Définition et déclaration des variables}
Faire une liste itemize des variables, avec leur domaine, ainsi que leur signification:

n blocs et k piles
Nous noterons actions(E) l'ensemble de toutes les actions possibles pour
un état E donné.
L'application d'une action i → j ∈ actions(E) sur un état E permet d'obtenir un nouvel état que nous noterons
t(E, i → j).
nous définissons le graphe d'états G = (S, A) tq:
— S est l'ensemble des états possibles, chaque état étant une configuration différente des n blocs sur les k piles ;
— A = {(E, E') | ∃i → j ∈ actions(E), E' = t(E, i → j)}, A est l'ensemble des arcs du graphe d'états, où t(E, i → j) est l'état obtenu à partir de l'état E en appliquant l'action i → j.

\subsection{Question 1}

\textit{Combien d'actions différentes sont-elles possibles pour l'état $E_1$ ?}

On a $k = 3$ piles non-vides à l'état $E_1$. Pour chaque pile, il y a $k - 1$ actions possibles : déplacer le bloc du sommet de la pile vers le sommet d'une \textit{autre} pile. On a donc $(k - 1) \times 3 = 6$ actions possibles pour l'état $E_1$.

\subsection{Question 2}
\textit{Étant donné un état $E$ de $n$ blocs sur $k$ piles, quelle est la taille maximale de $actions(E)$ ?}

On doit considérer le cas dans lequel il y a le plus d'actions possibles. Comme il n'y a pas d'action depuis une pile vide, on déduit que le nombre maximal d'action possible a lieu quand toutes les piles sont non-vides. On a donc $k$ piles non-vides, et pour chaque pile, on a $k - 1$ actions possibles. On a donc $k \times (k - 1)$ actions possibles.

\subsection{Question 3}

\textit{Étant donné un état $E$ de $n$ blocs sur $k$ piles ayant $v$ piles vides, quelle est la taille de $actions(E)$ ?}

reponse a developper: (k-v)(k-1)

\subsection{Question 4}

\textit{Quel est l'ordre de grandeur du nombre total d'états différents possibles ?}

reponse a développer: exponentiel

\section{Partie 2 : Définition du graphe d'états}

\subsection{Question 5}

\textit{Le graphe d'états $G$ est-il orienté ?}

Oui, orienté

\subsection{Question 6}

\textit{Quels sont les algorithmes qui peuvent être utilisés pour rechercher ce plus court chemin ?}

Bellman-Ford, BFS, Dijkstra

\subsection{Question 7}

\textit{Quel est l'algorithme le plus efficace pour rechercher ce plus court chemin ?}

BFS car le graph est non pondéré, donc pas besoin d'algorithmes plus complexes

\subsection{Question 8}

\textit{Quelle est la complexité en temps de cet algorithme par rapport à |S| et |A| ?}

linéaire

\subsection{Question 9}

\textit{Quelle est la complexité en temps de cet algorithme par rapport au nombre de blocs (n) et de piles (k) ?}

exponentielle

\section{Partie 3 : Heuristiques pour le monde des blocs}

\subsection{Question 10}

On commence par exécuter le programme de recherche de plus court chemin pour 4 piles, et un nombre de blocs variant de 6 à 8. On obtient les résultats suivants :

\begin{minipage}{\dimexpr\linewidth-20pt}
    \begin{lstlisting}[language=bash, caption={Résultat de l'algorithme PageRank après convergence sur le graphe $G_2$, avec $\alpha = 0.9$ et $\epsilon = 1.0 \times 10^{-10}$.}]
        $ make run
        Enter the number of stacks: 4
        Enter the number of blocs: 6
        Optimal solution of length 7 found in 12982 iterations and 0.060902 seconds
        [...]
        $ make run
        Enter the number of stacks: 4
        Enter the number of blocs: 7
        Optimal solution of length 9 found in 188569 iterations and 1.04769 seconds
        [...]
        $ make run
        Enter the number of stacks: 4
        Enter the number of blocs: 8
        Optimal solution of length 11 found in 2224481 iterations and 15.6446 seconds
        [...]
    \end{lstlisting}
\end{minipage}

On peut observer que le nombre d'itérations et le temps d'exécution augmentent de manière exponentielle avec le nombre de blocs. Pour améliorer les performances de l'algorithme, on peut utiliser une heuristique pour guider la recherche.

On introduit les heuristiques suivantes :
— h1 : Nombre de blocs ne se trouvant pas sur la dernière pile.
— h2 : Nombre de blocs ne se trouvant pas sur la dernière pile, plus deux fois le nombre de blocs b tels que b se
trouve sur la dernière pile mais il devra nécessairement être enlevé de cette pile pour ajouter et/ou supprimer
d'autres blocs sous lui.
— h3 : Nombre de blocs de Ei ne se trouvant pas sur la dernière pile, plus le nombre de blocs se trouvant au
dessus de chaque bloc ne se trouvant pas sur la dernière pile.

On dit qu'une heuristique est admissible si...
Let δ (u, t) be the distance between u and t, i.e. The length of a shortest path. h is admissible if foreach node u, h(u) ≤ δ (u, t).

une heuristique h est plus informée qu'une heuristique h' si pour tout état E, h(E) ≥ h'(E), tandis que les deux heuristiques sont
incomparables s'il existe deux états E et E' tels que h(E) < h'(E) et h(E') > h'(E')

\subsection{Question 11}

\textit{L'heuristique h1 est-elle admissible ?}

h1 < δ, donc oui

\subsection{Question 12}

h2 < δ, donc oui

\textit{L'heuristique h2 est-elle admissible ?}

\subsection{Question 13}

\textit{L'heuristique h2 est elle plus informée, moins informée, ou incomparable par rapport à l'heuristique h1 ?}

h1 < h2 < δ, donc h2 plus informée que h1

\subsection{Question 14}

\textit{L'heuristique h3 est-elle admissible ?}

h3 < δ, donc oui

\subsection{Question 15}

\textit{L'heuristique h3 est elle plus informée, moins informée, ou incomparable par rapport à l'heuristique h2 ?}

h3 < h2 < δ donc h3 moins informée que h2




%%%%%%%%%%%%%%%%%%%%%%%%%%%%%%%%%%%%%%%%%%%%%%%%%%%%%%%%%%%%%%%%%%%%%%%%%%%%%%%%%%%%%%%%%

\newpage
% glossary and acronyms
% \printglossary[type=\acronymtype, title=Acronymes]
% \printglossary[title=Glossaire]
\printglossary[type=\acronymtype]

\printglossary

% biblio
\printbibliography[
    heading=bibintoc,
    category=cited,
    title={Références}
]

% uncited references (bibliography)
% https://tex.stackexchange.com/questions/6967/how-to-split-bibliography-into-works-cited-and-works-not-cited
\printbibliography[
    notcategory=cited,
    heading=bibintoc,
    title={Bibliographie additionnelle},
]

\end{document}